\section{Domino tilings in three dimensions}\label{sec:three-dimensions}
We extend the previous tiling results for regular regions to three dimensional regions.
For integers $n_1, n_2, n_3 \geq 2$, define 
\begin{equation}
  G(n_1, n_2, n_3) = \Big\{(x_1, x_2, x_3) \in \mathbb{Z}^3 \Bigm| 0 \leq x_i \leq n_i \text{ for } i = 1,2,3 \Big\}
\end{equation}
to be the \emph{grid graph} with dimensions $n_1, n_2, n_3$. The grid graph can be divided into various regions based on how many coordinates are extremal (equal to $0$ or $n_i$). Specifically, for a grid graph $G \subseteq \mathbb{Z}^3$, define
\begin{equation}
  G^{[k]} = \Big\{p \in G \mid \text{exactly $3 - k$ coordinates of $p$ are extremal}\Big\}.
\end{equation}
For $k = 0, 1, 2, 3$, these are called \emph{corner}, \emph{edge}, \emph{face}, and \emph{interior} points, respectively.
In this setting, there are two domino configurations we wish to study: the familiar \emph{pair}, consisting of two dominoes sharing a long face as shown in Figure \ref{fig:3d-pair}, which we redefine here as
\begin{equation}
  D_{\mathrm{pair}} = \Big\{
      \big\{(0, 0, 0), (1, 0, 0)\big\}, 
      \big\{(0, 1, 0), (1, 1, 0)\big\}
  \Big\}
\end{equation}
and a new configuration, the \emph{loop}, consisting of three dominoes connected \say{head-to-tail} cyclically, as shown in Figure \ref{fig:3d-loop}. Formally, this can be written as
\begin{equation}
  D_{\mathrm{loop}} = \Big\{
    \big\{(0, 0, 0), (1, 0, 0)\big\},
    \big\{(0, 1, 0), (0, 1, 1)\big\},
    \big\{(1, 0, 1), (1, 1, 1)\big\}
  \Big\}.
\end{equation}

Three-dimensional tilings with pairs and loops have been previously studied by C.~Klivans and N.~Saldanha in \cite{klivans2025}.
For the loop substructure, they define an operation called the \emph{trit}, which modifies a tiling $T$ by replacing a loop with a copy of its chiral twin. 

\tdplotsetmaincoords{70}{30}
\begin{figure}[h]
  \centering
  \begin{minipage}{0.45\textwidth}
    \centering
    \begin{tikzpicture}[
      tdplot_main_coords,
      semithick,
      line join=round,
      facea/.style={fill=black!50},
      faceb/.style={fill=black!30},
      facec/.style={fill=black!40}
    ]
      % draw first domino
        % 1x1s
        %\draw (0, 0, -1) -- (1, 0, -1) -- (1, 1, -1) -- (0, 1, -1) -- cycle;
        \draw[facec] (0, 0, +1) -- (1, 0, +1) -- (1, 1, +1) -- (0, 1, +1) -- cycle;
      
        % 2x1s
        %\draw (0, 0, -1) -- (0, 0, +1) -- (0, 1, +1) -- (0, 1, -1) -- cycle;
        \draw[faceb] (0, 0, -1) -- (0, 0, +1) -- (1, 0, +1) -- (1, 0, -1) -- cycle;
        %\draw (1, 1, -1) -- (1, 1, +1) -- (0, 1, +1) -- (0, 1, -1) -- cycle;
        %\draw (1, 1, -1) -- (1, 1, +1) -- (1, 0, +1) -- (1, 0, -1) -- cycle;

      % draw second domino
      \begin{scope}[shift={(1,0,0)}]
        % 1x1s
        %\draw (0, 0, -1) -- (1, 0, -1) -- (1, 1, -1) -- (0, 1, -1) -- cycle;
        \draw[facec] (0, 0, +1) -- (1, 0, +1) -- (1, 1, +1) -- (0, 1, +1) -- cycle;
      
        % 2x1s
        %\draw (0, 0, -1) -- (0, 0, +1) -- (0, 1, +1) -- (0, 1, -1) -- cycle;
        \draw[faceb] (0, 0, -1) -- (0, 0, +1) -- (1, 0, +1) -- (1, 0, -1) -- cycle;
        %\draw (1, 1, -1) -- (1, 1, +1) -- (0, 1, +1) -- (0, 1, -1) -- cycle;
        \draw[facea] (1, 1, -1) -- (1, 1, +1) -- (1, 0, +1) -- (1, 0, -1) -- cycle;
      \end{scope}
    \end{tikzpicture}
    \caption{The pair configuration.}\label{fig:3d-pair}
  \end{minipage}
  \begin{minipage}{0.45\textwidth}
    \centering
    \begin{tikzpicture}[
      tdplot_main_coords,
      semithick,
      line join=round,
      facea/.style={fill=black!50},
      faceb/.style={fill=black!30},
      facec/.style={fill=black!40}
    ]
      % draw third domino
      \begin{scope}[rotate around y=-90, shift={(-1,1,-1)}]
        % 1x1s
        \draw[facea] (0, 0, -1) -- (1, 0, -1) -- (1, 1, -1) -- (0, 1, -1) -- cycle;
        %\draw (0, 0, +1) -- (1, 0, +1) -- (1, 1, +1) -- (0, 1, +1) -- cycle;
      
        % 2x1s
        %\draw (0, 0, -1) -- (0, 0, +1) -- (0, 1, +1) -- (0, 1, -1) -- cycle;
        \draw[faceb] (0, 0, -1) -- (0, 0, +1) -- (1, 0, +1) -- (1, 0, -1) -- cycle;
        %\draw (1, 1, -1) -- (1, 1, +1) -- (0, 1, +1) -- (0, 1, -1) -- cycle;
        %\draw (1, 1, -1) -- (1, 1, +1) -- (1, 0, +1) -- (1, 0, -1) -- cycle;
      \end{scope}

      % draw first domino
        % 1x1s
        %\draw (0, 0, -1) -- (1, 0, -1) -- (1, 1, -1) -- (0, 1, -1) -- cycle;
        \draw[facec] (0, 0, +1) -- (1, 0, +1) -- (1, 1, +1) -- (0, 1, +1) -- cycle;
      
        % 2x1s
        %\draw (0, 0, -1) -- (0, 0, +1) -- (0, 1, +1) -- (0, 1, -1) -- cycle;
        \draw[faceb] (0, 0, -1) -- (0, 0, +1) -- (1, 0, +1) -- (1, 0, -1) -- cycle;
        %\draw (1, 1, -1) -- (1, 1, +1) -- (0, 1, +1) -- (0, 1, -1) -- cycle;
        \draw[facea] (1, 1, -1) -- (1, 1, +1) -- (1, 0, +1) -- (1, 0, -1) -- cycle;

      % draw second domino
      \begin{scope}[shift={(1,1,0)}, rotate around x=90]
        % 1x1s
        %\draw (0, 0, -1) -- (1, 0, -1) -- (1, 1, -1) -- (0, 1, -1) -- cycle;
        \draw[faceb] (0, 0, +1) -- (1, 0, +1) -- (1, 1, +1) -- (0, 1, +1) -- cycle;
      
        % 2x1s
        %\draw (0, 0, -1) -- (0, 0, +1) -- (0, 1, +1) -- (0, 1, -1) -- cycle;
        %\draw (0, 0, -1) -- (0, 0, +1) -- (1, 0, +1) -- (1, 0, -1) -- cycle;
        \draw[facec] (1, 1, -1) -- (1, 1, +1) -- (0, 1, +1) -- (0, 1, -1) -- cycle;
        \draw[facea] (1, 1, -1) -- (1, 1, +1) -- (1, 0, +1) -- (1, 0, -1) -- cycle;
      \end{scope}
    \end{tikzpicture}
    \caption{The loop configuration.}\label{fig:3d-loop}
  \end{minipage}
\end{figure}

\noindent While it has been proven in \cite{thurston1990} that tilings of simply connected regions of the plane are flip-connected, it is not yet known whether or not even the tilings of $G(n_1, n_2, n_3)$ are connected by flips and trits.
Since every grid graph has more than one domino tiling, a necessary condition for this to be true is for each tiling to contain at least one pair or loop.
We show that this is the case.


\begin{theorem}\label{thm:grid}
  Let $n_1, n_2, n_3 \geq 2$ be integers.
  If $T$ is a domino tiling of $G(n_1, n_2, n_3)$, then $T$ contains a pair or a loop.
\end{theorem}
\begin{proof}
  Assume, for the sake of contradiction, that there exists a tiling $T$ which does not contain any substructure isomorphic to $D_{\mathrm{pair}}$ or $D_{\mathrm{loop}}$.
  We first show that every subcube $Q$ must contain at most two whole dominoes.
  By a suitable translation, we may assume that $Q = \{0,1\}^3$. 
  Suppose there are already at least two dominoes fully contained within $Q$, and that they do not already form a pair.
  Then there are two cases:
  \begin{enumerate}
    \item If the two dominos are parallel, then by a suitable isometry we may take these to be
    \begin{equation}
      \{\delta_1, \delta_2\} = \Big\{
        \big\{(0, 0, 0), (0, 0, 1)\big\}, 
        \big\{(1, 1, 0), (1, 1, 1)\big\}
      \Big\}.
    \end{equation}
    There are two remaining candidates for a possible third domino $\delta_3 \subseteq Q$, 
    but, as one can verify, both produce a substructure isomorphic to $D_{\mathrm{pair}}$.
    
    \item Otherwise, by a suitable isometry we may take these to be
    \begin{equation}
      \{\delta_1, \delta_2\} = \Big\{
        \big\{(0, 0, 0), (1, 0, 0)\big\}, 
        \big\{(0, 1, 0), (1, 1, 0)\big\}
      \Big\}.
    \end{equation}
    There are three remaining candidates for a possible third domino $\delta_3 \subseteq Q$, 
    but, as one can verify, two produce a substructure isomorphic to $D_{\mathrm{pair}}$, 
    and the other yields a substructure isomorphic to $D_{\mathrm{loop}}$.
  \end{enumerate}
  It follows that
  \begin{equation}\label{eq:3d-local-div}
    (\operatorname{div} T)(Q) 
    = 1 - \tfrac{1}{4} \cdot \big|\{q \in Q \mid T(q) \subseteq Q \}\big|
    \geq 1 - \tfrac{1}{4} \cdot 2 \cdot 2
    = 0. 
  \end{equation}
  for all subcubes $Q$.
  For the rest of the argument, we calculate the sum of $\operatorname{div} T$ over each cube in the region, reduce it to a sum over the boundary of the grid.

  Let $q \in G^{[k]}$. 
  If $k > 0$, then there exist $2^k$ subcubes containing $q$, and at least half of these cubes will fully contain $\delta_T(q)$.
  Hence,
  \begin{equation}
    \sum_{\substack{Q \mskip 1mu \subseteq \mskip 0.3mu G \\[0.2ex] Q \ni q}} \iota(\delta_T(q); Q) \leq 2^{k-1} - 2^{k-1} = 0.
  \end{equation}
  Otherwise, if $k = 0$, then the one subcube containing $q$ also fully contains $\delta_T(q)$, so 
  \begin{equation}
    \sum_{\substack{Q \mskip 1mu \subseteq \mskip 0.3mu G \\[0.2ex] Q \ni q}} \iota(\delta_T(q); Q) = -1.
  \end{equation}
  Overall, we obtain
  \begin{equation}
    \sum_{Q \subseteq G} (\operatorname{div} T)(Q) =
    \frac{1}{8} \cdot \sum_{q \in G} \sum_{Q \ni q} f(T(q), Q) 
    \leq -1,
  \end{equation}
  which contradicts \eqref{eq:3d-local-div}. 
  Therefore, every domino tiling $T$ of the grid graph $G(n_1, n_2, n_3)$ must contain a substructure isomorphic to $D_{\mathrm{pair}}$ or $D_{\mathrm{loop}}$.
\end{proof}